Satellite imaging of tropical rainforests is limited by persistent cloud cover.
Conventional controllers image every target on a fixed pointing schedule, regardless of cloud conditions.
Shutter budget and downlink capacity are wasted on cloud-obscured frames scientists cannot use.
An autonomous pointing and scheduling policy that reacts to onboard camera observations would use those resources more efficiently.

Because that idea appears straightforward, this work asks a narrower question:
what reward shaping and action-space design are required to train such a policy in simulation?
We build a reproducible platform for that investigation: single-axis attitude dynamics, ray-tracing cameras, parameterised cloud primitives, and a multi-phase safety controller.
A deterministic sequential-target overflight policy serves as the comparison baseline.
Two off-policy reinforcement-learning algorithms are trained and evaluated: Soft Actor-Critic and Maximum a Posteriori Policy Optimisation.

The experiments show that the problem is not trivial in practice.
Reward structure, how the agent commands attitude, and correct algorithm implementation each had to be resolved before useful behaviour emerged.
Both algorithms eventually develop selective shuttering absent in the deterministic baseline, but reliable mission-level gains over that baseline were not established within the semester scope.

The study is limited to a 2D single-pass orbital model.
Attitude safety constraints are enforced throughout.
Flight hardware, multi-orbit planning, and three-dimensional operations are out of scope.
